\small
\begin{tabularx}{\linewidth}{@{}>{\raggedright\arraybackslash}p{2.6cm}>{\raggedright\arraybackslash}X@{}}
\toprule
\textbf{Quem} & \textbf{Slides, na ordem}\\
\midrule
\textbf{Kauan} & Capa; Roteiro; 1 · Questão de pesquisa; Limitações dos dados; Escopo: municípios com \nPopMin{} mil hab. ou mais\\
\textbf{Fabio} & 2 · Base de dados; 3 · Variáveis\\
\textbf{Douglas} & 4 · A taxa é assimétrica; o log corrige; 4 · A violência se concentra no Norte e Nordeste; 4 · Nenhum fator sozinho explica a taxa\\
\textbf{Luiz Felipe} & 4 · A renda tem efeito em U invertido; 5 · Box-Cox indica o logaritmo da taxa; 5 · Construção do modelo em blocos\\
\textbf{Willian} & 5 · Modelo preliminar: efeito de cada variável na taxa; 5 · Efeito da renda previsto pelo modelo; 5 · Diagnóstico dos resíduos; Próximos passos; Obrigado\\
\bottomrule
\end{tabularx}
\normalsize
